\section{Setup and Preliminaries}\label{sec:setting}

\paragraph{Notations.} We use $\|\cdot\|$ to denote the Euclidean 2-norm. For a symmetric positive definite matrix $A$ and a vector $x$ of proper dimension, we let $\|x\|_{A}\coloneqq \sqrt{x^\top A x}$ be the weighted 2-norm (or the $A$-norm). We let $\dotp{\cdot}{\cdot}$ denote the inner product in Euclidean  space such that $\dotp{x}{y} = x^\top y$. For a $d$-dimensional matrix $V$,   we let $\lambda_1(V)\geq  \lambda_2(V) \geq  \dots \geq  \lambda_d(V)$ be the eigenvalues of $V$  arranged in decreasing order. 
We let $\cB_d $ denote the unit ball in $\RR^d$, and $\cS_{d-1} = \{x\in\RR^d:\|x\| =  1\}$ denote the unit hypersphere in  $\RR^d$. For an interger $N\geq 1$, we let $[N]$ denote the set $\{1,2,\dots,N\}$. We use the $\cO(\cdot	)$ notation to suppress problem-dependent constants, and the $\tilde{\cO}(\cdot)$ notation further suppresses polylog factors.

\paragraph{Problem formulation and assumptions.} We consider the stochastic linear bandit problem. Let $\theta^\star \in \RR^d$ be a fixed but unknown parameter. At each time $t\in[T]$, a policy $\pi$  selects action $x_t$ from a set of action  $\cX_t \subset \RR^d $ according to the past observations and receives a reward $r_t  = \dotp{x_t}{\theta^\star} + \varepsilon_t$, where $ \varepsilon_t$ is mean-zero noise with a distribution specified in Assumption \ref{assump:noise} below. We measure the performance of $\pi$ with the cumulative expected regret 
%
$
\cR(T) = \sum_{t= 1}^T\dotp{x^{\star}_t}{\theta^\star} -  \dotp{x_t}{\theta^\star}\,,
$
%
where $x^{\star}_t$ is the best action at time $t$, i.e., 
%
$
x^{\star}_t=\arg\max_{x \in \cX_t}\dotp{x}{\theta^\star}\,.
$
%
Let $\mathcal{F}_t$ be a $\sigma$-algebra generated by the history $(x_1, r_1, \ldots, x_t, r_t)$ and the prior knowledge, $\mathcal{F}_0$. Therefore, $\{\mathcal{F}_t\}_{t\ge 0}$ forms a filteration such that each $\mathcal{F}_t$ encodes all the information up to the end of period $t$.

We make the following assumptions that are standard in the relevant literature.

\begin{assumption}[Bounded parameter]\label{assump:bounded-parameter}
	The unknown parameter  $\param$ is bounded as  $\|\param\|\leq S$, where $S > 0$ is known.
\end{assumption}
\begin{assumption}[Bounded action sets]\label{assump:bounded-action-sets}
	The action sets $\{\cX_t\}$ are uniformly bounded and closed subsets of $\RR^d$, such that $\|x\|\leq X_t$ for all $x\in \cX_t$ and all $t\in[T]$, where $X_t$'s are known and $\sup_{t\geq 1}\left\{X_t\right\} < \infty $. %Without loss of generality, we normalize the actions such that $\cX_t \subset \cS_{d-1}$ for all $t\in[T]$.
\end{assumption}


\begin{assumption} [Subgaussian reward noise]\label{assump:noise}
	The noise sequence $\{\varepsilon_t\}_{t\ge 1}$ is conditionally mean-zero and $R$-subgaussian, where $R$ is known. Formally, 
	for all real valued $\lambda$, 
	$
	 \EE\left[e^{\lambda \varepsilon_t} | \cF_t\right] \leq \exp \left(\lambda^2 R^2 / 2\right)
	 $.
	 This condition implies that $\EE \left[\varepsilon_t  | \cF_t\right]=0$ for all $t\geq 1$.
\end{assumption}



\subsection{Regularized  Least Square and Confidence Ellipsoid}

In this subsection, we review the useful frequentist tools developed by \cite{abbasi2011improved} for estimating the unknown parameter $\theta^*$ in linear bandit (LB) problems.

Consider an arbitrary sequence of actions $(x_1,\dots,x_t)$ and their corresponding rewards $(r_1,\dots,r_t)$. In LB problems, the parameter $\theta^*$ is typically estimated using the regularized least squares (RLS) estimator. Let $\lambdareg$ be a fixed regularization parameter. The sample covariance matrix $V_t$ and the RLS estimate $\hat{\theta}_t$ are defined as follows:
\begin{equation}\label{eq:ridge}
	V_t=\lambdareg I_d +\sum_{s=1}^{t} x_s x_s^{\top}, \quad \hat{\theta}_t=V_t^{-1} \sum_{s=1}^{t} x_s r_{s}.
\end{equation}

The following proposition from \cite{abbasi2011improved} establishes that the RLS estimate $\hat{\theta}_t$ concentrates around the true parameter $\theta^*$ with high probability.

\begin{proposition}[Theorem 2 in \cite{abbasi2011improved}]\label{prop:beta}
	Let $\delta \in (0,1)$ be a fixed confidence level. Then, with probability at least $1-\delta$, it holds for all $x \in \mathbb{R}^d$ that
	\begin{align*}
		\|\widehat{\theta}_t-\theta^{\star}\|_{V_t} \leq \beta_{t,\delta, \lambdareg}^{RLS}, \quad 
		|\dotp{x}{\widehat{\theta}_t-\theta^{\star}}| \leq\|x\|_{V_t^{-1}} \beta_{t,\delta, \lambdareg}^{RLS}
	\end{align*}
	where the confidence bound $\beta_{t,\delta, \lambdareg}^{RLS}$ is defined as
	\begin{equation}\label{eq:beta_rls}
		\beta_{t,\delta, \lambdareg}^{RLS} = R \sqrt{2 \log (\lambdareg+t)^{d / 2} \lambdareg^{-d / 2}{\delta}^{-1}} + \sqrt{\lambdareg} S.
	\end{equation}
\end{proposition}

 Proposition~\ref{prop:beta} enables us to construct the following sequence of confidence ellipsoids.
\begin{definition}
	\label{def:rls}
	Fix $\delta\in (0,1)$. We define the RLS confidence ellipsoid as	
	\[
	\cE_{t,\delta, \lambdareg}^{RLS}= \{\theta \in \RR^d: \|\theta - \hparam_t\|_{V_t}\leq \beta_{t,\delta, \lambdareg}^{RLS} \}\,.
	\]
\end{definition}

The next proposition, known as the \emph{elliptical potential lemma}, plays a central role in bounding the regret. This proposition provides the key element in the work of \cite{abbasi2011improved}, showing that the cumulative prediction error incurred by the action sequence used to estimate $\theta^*$ is small.

\begin{proposition}[Lemma 11 in \cite{abbasi2011improved}] \label{prop:potential}
	If $\lambdareg >1$, for an arbitrary sequence $(x_1,\dots,x_t)$, it holds that
	$
	\sum_{s=1}^t\left\|x_s\right\|_{V_s^{-1}}^2 \leq 2 \log \frac{\operatorname{det}\left(V_{t+1}\right)}{\operatorname{det}(\lambdareg I)} \leq 2 d \log (1+\frac{t}{\lambdareg})\,.
	$

\end{proposition}


\section{\POFUL{} Algorithms}
\label{sec:POFUL}
In this section, we introduce \POFUL{} (Pivot OFUL), a generalized framework of OFUL. This framework enables a unified analysis of frequentist regret for common algorithms.

At a high level, \POFUL{} is designed to encompass the exploration mechanism of OFUL and LinTS. \POFUL{} takes as input a sequence of \emph{inflation} parameters $\{\iota_t\}_{t\in [T]}$,  feasible (randomized) \emph{pivots} $\{\tilde{\theta}_t\}_{t\in[T]}$ and \emph{optimism} parameters $\{\tau_t\}_{t\in [T]}$. The inflation parameters are used to construct confidence ellipsoids that contain $\{\tilde{\theta}_t\}_{t\in[T]}$ with high probability. This is formalized in the next definition.%Formally,  $\{\iota_t\}_{t\in [T]}$ and $\{\tilde{\theta}_t\}_{t\in[T]}$ 

\begin{definition} \label{def:tilde_beta}
Fix  $\delta \in (0,1)$ and  $\delta^\prime = \delta/2 T$. Given the inflation parameters $\{\iota_t\}_{t\in [T]}$, we call random variables $\{\tilde{\theta}_t\}_{t\in[T]}$  feasible pivots if for all $t\in[T]$,
%
$
	\PP[\tilde{\theta}_t \in \cE_{t,\delta^\prime, \lambdareg}^{\PVT} |\mathcal{F}_t] \geq 1 - \delta^\prime
$,
% 
where we define the ``pivot ellipsoid'' as
$
 \cE_{t,\delta, \lambdareg}^{\PVT}  \coloneqq  \{\theta \in \RR^d: \|\theta - \hparam_t\|_{V_t}\leq \iota_t \beta^{RLS}_{t,\delta,\lambdareg} \}
$.
\end{definition}

At each time $t$, \POFUL{} chooses the action that maximizes the optimistic reward 
%
\#\label{eq:POFUL}
\tilde{x}_t = \argmax_{x\in\cX_t} \  \dotp{x}{\tilde{\theta}_{t}}  +   \tau_t  \|x\|_{V^{-1}_t}\beta_{t,\delta^\prime, \lambdareg}^{RLS}\,,
\#
%
as shown in a pseudocode representation in Algorithm~\ref{alg:POFUL} and illustrated in Figure~\ref{fig:POFUL_illustration}. 

Recall OFUL encourages exploration by introducing the uncertainty term $\tau_t  \|x\|_{V^{-1}_t}\beta_{t,\delta^\prime, \lambdareg}^{RLS}$ in the reward, while LinTS explores through random sampling within the confidence ellipsoid. We let \POFUL{} select an arbitrary pivot (which can be random) from $\cE_{t,\delta^\prime, \lambdareg}^{\PVT}$ and maximize the optimistic reward to encompass arbitrary exploration mechanisms within $\cE_{t,\delta^\prime, \lambdareg}^{\PVT}$. 

\begin{algorithm}
	\caption{\POFUL{} }\label{alg:POFUL}
	\begin{algorithmic}
		% \Require $n \geq 0$
		% \Ensure $y = x^n$
		\REQUIRE $T$, $\delta$, $\lambdareg$,  $\{\iota_t\}_{t\in [T]}$, $\{\tau_t\}_{t\in [T]}$
		\STATE Initialize $ V_{0} \leftarrow \lambdareg I_d$, $\widehat{\theta}_{1} \leftarrow 0$,  $\delta^\prime \leftarrow \delta/2T$
		\FOR{t = 0, 1,  \dots, T }
		\STATE Sample a feasible pivot $\tilde{\theta}_{t} $ with respect to  $\iota_t$ according to Definition~\ref{def:tilde_beta}
		% \STATE Sample $\tilde{\theta}_{t} \sim \mathcal{N}\left(\hat{\theta}_{t}, \iota^{2}_t  V_{t}^{-1}\right)$
		\STATE $\tilde{x}_{t} \leftarrow \argmax_{x\in\cX_t} \dotp{x}{\tilde{\theta}_{t}}  +   \tau_t  \|x\|_{V^{-1}_t}\beta_{t,\delta^\prime, \lambdareg}^{RLS}$
		\STATE Observe reward $r_{t}$
		\STATE $V_{t+1} \leftarrow V_{t}+\tilde{x}_{t}  \tilde{x}_{t}^{\top}$
		\STATE $\hat{\theta}_{t+1} \leftarrow V_{t+1}^{-1}\sum_{s=1}^{t} \tilde{x}_s r_{s}$.
		\ENDFOR
	\end{algorithmic}
\end{algorithm}


\begin{figure}[!b]
\centering
\begin{subfigure}[]{0.48\textwidth}
	\centering
 
	\begin{tikzpicture}[scale=0.5, every node/.append style={transform shape, font=\Large}]
		
		% Adjust the position of the left-half figure
		\def\originX{-8}
		\def\originY{0}
		
		% Define the centers and radiuses of ellipsoids
		\def\centerxhat{\originX+4}
		\def\centeryhat{\originY+4}
		
		\def\centerxtilde{\centerxhat-1}
		\def\centerytilde{\centeryhat+1.5}
		
		\def\radiusHatX{3} 
		\def\radiusHatY{1.5}
		\def\radiusTildeX{1.5} 
		\def\radiusTildeY{0.75}
		
		%% Draw the ellipsoid
		
		% Draw axises 
		\filldraw [black] (\originX,\originY) circle (1pt) node[anchor=north east] {$\mathcal{O}$};
		\draw[->,axiscolor] (\originX,\originY+0) -- (\originX+8,\originY+0) node[below] {};
		\draw[->,axiscolor] (\originX+0,\originY) -- (\originX+0,\originY+7.5) node[left] {};
		
		% Draw the PVT ellipsoid
		\draw [rotate around = {135:(\centerxhat,\centeryhat)}] (\centerxhat,\centeryhat) ellipse (\radiusHatX cm and \radiusHatY cm);
		\filldraw [] (\centerxhat,\centeryhat) circle (0.5pt) node[inner sep = 0pt, anchor=south west, font=\Large] {$\hat{\theta}_t$}; 
		\node[] at (\centerxhat+3.2,\centeryhat-0.7) { \Large $\mathcal{E}_{t,\delta^\prime, \lambdareg}^{P V T}$};
		
		% Draw the OFUL ellipsoid
		\draw [rotate around = {135:(\centerxtilde,\centerytilde)}, pastelblue] (\centerxtilde,\centerytilde) ellipse (\radiusTildeX cm and \radiusTildeY cm);
		\filldraw [] (\centerxtilde,\centerytilde) circle (0.5pt) node[inner sep = 0pt, anchor=south west, font=\Large] {$\tilde{\theta}_t$}; 
		
		
		% Vertical line to split the figure
		\draw[line width=2pt] (\originX+9,\originY) -- (\originX+9,\originY+7.5);
		
		% Add text for "parameter space"
		\node[text width=4cm, align=center, font=\Large] at (\originX+4, \originY+.5) {Parameter space: $\mathbb{R}^d$};
		
		
		
		%% Draw the intervals 
		
		% Adjust the position of the right-half figure
		\def\intervalAxisX{\centerxhat+8}
		\def\intervalAxisY{\centeryhat} % matching the centers
		
		% Add text for "reward space"
		\node[text width=3cm, align=center, font=\Large] at (\originX+12, \originY+.5) {Reward space: $\mathbb{R}$};
		
		
		\draw[->,axiscolor] (\intervalAxisX,\originY+1) -- (\intervalAxisX,\originY+7.5) node[left] {};
		
		% Draw short horizontal lines (segments)
		
		% Calculate the projected radius for pi/4 clockwise rotation 
		\pgfmathsetmacro{\radiusPVT}{(\radiusHatX)^2*sqrt(2)/2/sqrt((\radiusHatY)^2 + (\radiusHatX)^2) + (\radiusHatY)^2*sqrt(2)/2/sqrt((\radiusHatY)^2 + (\radiusHatX)^2)}
		
		\pgfmathsetmacro{\radiusOFUL}{(\radiusTildeX)^2*sqrt(2)/2/sqrt((\radiusTildeY)^2 + (\radiusTildeX)^2) + (\radiusTildeY)^2*sqrt(2)/2/sqrt((\radiusTildeY)^2 + (\radiusTildeX)^2)}
		

		%  Draw POFUL parts
		\node[anchor=west, font=\Large] (POFUL) at (\intervalAxisX+2,\centerytilde+\radiusOFUL) {$\text{\POFUL}$};
		\draw[dashed] (POFUL) -- (\intervalAxisX+0,\centerytilde+\radiusOFUL); 
		
		\draw[decorate,decoration={brace, amplitude=10pt, aspect = 0.3, mirror}] (\intervalAxisX,\centeryhat-\radiusPVT) -- (\intervalAxisX,\centeryhat+\radiusPVT) node[pos = 0.3, xshift=20pt, anchor=west, font=\Large] { 2$\iota_t\|x\|_{V_t^{-1}} \beta_{t,\delta^\prime,\lambdareg}^{R L S}$};
		
		\filldraw [] (\intervalAxisX,\centerytilde) circle (1pt) node[inner sep = 2.5pt, anchor=  east, font=\Large] {$\langle\tilde{\theta}_t, x\rangle$};
		
		
		
		\draw[decorate,decoration={brace, amplitude=5pt, mirror,aspect=0.3},pastelblue] (\intervalAxisX,\centerytilde - \radiusOFUL) -- (\intervalAxisX, \centerytilde+\radiusOFUL) node[pos = 0.3, xshift= 10pt, anchor=west, font=\Large] { 2$\tau_t\|x\|_{V_t^{-1}} \beta_{t,\delta^\prime,\lambdareg}^{R L S}$};
		
		\filldraw [black] (\intervalAxisX,\intervalAxisY) circle (1pt) node[inner sep = 2.5pt, anchor= east, font=\Large] {$\langle\hat{\theta}_t, x\rangle$};
		
		\draw (\intervalAxisX-0.1,\centeryhat+\radiusPVT) -- (\intervalAxisX+0.1,\centeryhat+\radiusPVT) node[right]{};
		\draw (\intervalAxisX-0.1,\centeryhat-\radiusPVT) -- (\intervalAxisX+0.1,\centeryhat-\radiusPVT) node[right]{};  
		\draw (\intervalAxisX-0.1,\centerytilde+\radiusOFUL) -- (\intervalAxisX+0.1,\centerytilde+\radiusOFUL)[pastelblue] node [right,pastelblue]{};
		\draw (\intervalAxisX-0.1,\centerytilde-\radiusOFUL) -- (\intervalAxisX+0.1,\centerytilde-\radiusOFUL)[pastelblue] node[right,pastelblue]{};
		
		
		% dashed lines connecting two half figures 
		\draw[dashed, dashedlinecolor] (\originX,\centerytilde+\radiusOFUL) -- (\intervalAxisX+0,\centerytilde+\radiusOFUL); 
		\draw[dashed, dashedlinecolor] (\originX,\centerytilde-\radiusOFUL) -- (\intervalAxisX+0,\centerytilde-\radiusOFUL); 
		
		\draw[dashed,dashedlinecolor] (\originX,\centeryhat+\radiusPVT) -- (\intervalAxisX+0,\centeryhat+\radiusPVT); 
		\draw[dashed,dashedlinecolor] (\originX,\centeryhat-\radiusPVT) -- (\intervalAxisX+0,\centeryhat-\radiusPVT);     
		
	\end{tikzpicture}
	\caption{}
	\label{fig:POFUL_illustration}
\end{subfigure}
\hfill
\begin{subfigure}[]{0.48\textwidth}
	\centering
	\begin{tikzpicture}[scale=0.5, every node/.append style={transform shape, font=\Large}]
		
		% Adjust the position of the left-half figure
		\def\originX{-8}
		\def\originY{0}
		
		% Define the centers and radiuses of ellipsoids
		\def\centerxhat{\originX+4}
		\def\centeryhat{\originY+4}

		\def\centerxtilde{\centerxhat-1}
		\def\centerytilde{\centeryhat+1.5}
		
		\def\radiusHatX{3} 
		\def\radiusHatY{1.5}
		\def\radiusTildeX{1.5} 
		\def\radiusTildeY{0.75}
		
		
		%% Draw the intervals 
		
		\def\intervalAxisX{\centerxhat-2}
		\def\intervalAxisY{\centeryhat} % matching the centers
		
		
		% Calculate the projected radius for pi/4 clockwise rotation 
		\pgfmathsetmacro{\radiusPVT}{(\radiusHatX)^2*sqrt(2)/2/sqrt((\radiusHatY)^2 + (\radiusHatX)^2) + (\radiusHatY)^2*sqrt(2)/2/sqrt((\radiusHatY)^2 + (\radiusHatX)^2)}
		
		\pgfmathsetmacro{\radiusOFUL}{(\radiusTildeX)^2*sqrt(2)/2/sqrt((\radiusTildeY)^2 + (\radiusTildeX)^2) + (\radiusTildeY)^2*sqrt(2)/2/sqrt((\radiusTildeY)^2 + (\radiusTildeX)^2)}

  

    \def\dodgeLength{.02}
    \def\PVTtickerlength{.09}
    \def\OFULtickerlength{.12}
    
% Draw OFUL parts
		
		\pgfmathsetmacro{\radiusOFUL}{0.6}
		
		\def\intervalAxisXOFUL{\intervalAxisX}
		
		\draw[->,axiscolor] (\intervalAxisXOFUL,\originY) -- (\intervalAxisXOFUL,\originY+7.5) node[left] {};
		\node[anchor=west, font=\Large] (OFUL) at (\intervalAxisXOFUL+1,\centeryhat+\radiusOFUL) {$\text{OFUL}$};
		\draw[dashed] (OFUL) -- (\intervalAxisXOFUL+0,\centeryhat+\radiusOFUL);
		%\draw[decorate,decoration={brace, amplitude=8pt, mirror,aspect=0.5},pastelblue] (\intervalAxisXOFUL,\centeryhat-\radiusOFUL) -- (\intervalAxisXOFUL,\centeryhat+ \radiusOFUL) node[pos=0.5, xshift= 10pt, anchor= west] { 2$\|x\|_{V_t^{-1}} \beta_t^{R L S}\left(\delta^{\prime}\right)$};
		
		
		\filldraw [black] (\intervalAxisXOFUL,\intervalAxisY) circle (1pt) node[inner sep = 2.5pt, anchor= east, font=\Large] {$\langle\hat{\theta}_t, x\rangle$}; 
		\draw (\intervalAxisXOFUL-\PVTtickerlength,\centeryhat+\dodgeLength) -- (\intervalAxisXOFUL+\PVTtickerlength,\centeryhat+\dodgeLength) node[right]{};
		\draw (\intervalAxisXOFUL-\PVTtickerlength,\centeryhat-\dodgeLength) -- (\intervalAxisXOFUL+\PVTtickerlength,\centeryhat-\dodgeLength) node[right]{};   
		\draw (\intervalAxisXOFUL-\OFULtickerlength,\centeryhat+\radiusOFUL) -- (\intervalAxisXOFUL+\OFULtickerlength,\centeryhat+\radiusOFUL)[pastelblue] node[right,pastelblue]{};
		\draw (\intervalAxisXOFUL-\OFULtickerlength,\centeryhat-\radiusOFUL) -- (\intervalAxisXOFUL+\OFULtickerlength,\centeryhat-\radiusOFUL)[blue] node[right,pastelblue]{};	
		
		% Draw TS parts
        \def\intervalAxisXTS{\intervalAxisX + 4}
		
		\draw[->,axiscolor] (\intervalAxisXTS,\originY) -- (\intervalAxisXTS,\originY+7.5) node[left] {};
		

		\node[anchor=west, font=\Large] (TS) at (\intervalAxisXTS+1,\centerytilde) {$\text{TS}$};
		\draw[dashed] (TS) -- (\intervalAxisXTS+0,\centerytilde); 
		
		\filldraw [] (\intervalAxisXTS,\centerytilde) circle (1pt) node[inner sep = 2.5pt, anchor=  east, font=\Large] {$\langle\tilde{\theta}_t, x\rangle$};
		
		\filldraw [black] (\intervalAxisXTS,\intervalAxisY) circle (1pt) node[inner sep = 2.5pt, anchor= east, font=\Large] {$\langle\hat{\theta}_t, x\rangle$};
		
		
		\draw (\intervalAxisXTS-\PVTtickerlength,\centeryhat+\radiusPVT) -- (\intervalAxisXTS+\PVTtickerlength,\centeryhat+\radiusPVT) node[right]{};
		\draw (\intervalAxisXTS-\PVTtickerlength,\centeryhat-\radiusPVT) -- (\intervalAxisXTS+\PVTtickerlength,\centeryhat-\radiusPVT) node[right]{};  
		\draw (\intervalAxisXTS-\OFULtickerlength,\centerytilde+\dodgeLength) -- (\intervalAxisXTS+\OFULtickerlength,\centerytilde+\dodgeLength)[pastelblue] node [right,pastelblue]{};
		\draw (\intervalAxisXTS-\OFULtickerlength,\centerytilde-\dodgeLength) -- (\intervalAxisXTS+\OFULtickerlength,\centerytilde-\dodgeLength)[pastelblue] node[right,pastelblue]{};
		
		
		
		% Draw Greedy parts
		
		\def\intervalAxisXGreedy{\intervalAxisX+8}
		
		\draw[->,axiscolor] (\intervalAxisXGreedy,\originY) -- (\intervalAxisXGreedy,\originY+7.5) node[left] {};
		
		\node[anchor=west, font=\Large] (Greedy) at (\intervalAxisXGreedy+1,\intervalAxisY) {$\text{Greedy}$};
		\draw[dashed] (Greedy) -- (\intervalAxisXGreedy+0,\intervalAxisY); 
		
		\filldraw [black] (\intervalAxisXGreedy,\intervalAxisY) circle (1pt) node[inner sep = 2.5pt, anchor= east, font=\Large] {$\langle\tilde{\theta}_t, x\rangle$=$\langle\hat{\theta}_t, x\rangle$};
		
		\draw (\intervalAxisXGreedy-\OFULtickerlength,\centeryhat+2*\dodgeLength) -- (\intervalAxisXGreedy+\OFULtickerlength,\centeryhat+2*\dodgeLength)[pastelblue] node [right,pastelblue]{};
		\draw (\intervalAxisXGreedy-\OFULtickerlength,\centeryhat-2*\dodgeLength) -- (\intervalAxisXGreedy+\OFULtickerlength,\centeryhat-2*\dodgeLength)[pastelblue] node[right,pastelblue]{};
		
		
		\draw (\intervalAxisXGreedy-\PVTtickerlength,\centeryhat+\dodgeLength) -- (\intervalAxisXGreedy+\PVTtickerlength,\centeryhat+\dodgeLength) node[right]{};
		\draw (\intervalAxisXGreedy-\PVTtickerlength,\centeryhat-\dodgeLength) -- (\intervalAxisXGreedy+\PVTtickerlength,\centeryhat-\dodgeLength) node[right]{};  
		
		
		
		
		
	\end{tikzpicture}
	\caption{}
	\label{fig:POFUL_special_cases}
 \end{subfigure}
 \caption{(a) POFUL algorithms illustration for general $\iota_t$ and  $\tau_t$. (b) Special cases: OFUL ($\iota_t = 0$, $\tau_t = 1$), TS ($\tau_t = 0$), and Greedy ($\iota_t = \tau_t = 0$).}
%\label{fig:POFUL_illustrations}
\end{figure}


We demonstrate that \POFUL{} encompasses OFUL, LinTS, TS-Freq, and Greedy as special cases, as illustrated in Figure \ref{fig:POFUL_special_cases}.
%
\begin{example}[OFUL]\label{ex:oful}
	For stochastic linear bandit problems, OFUL chooses actions by solving the optimization problem 
	$
	\max_{x\in\cX_t}  \   \dotp{x}{ \hat{\theta}_t}  + \|x\|_{V_t^{-1}}   \beta_{t,\delta^\prime, \lambdareg}^{RLS}
	$.
	%	
	Therefore, OFUL is a specially case of \POFUL{} where $\iota_t = 0$, $\tau_t = 1$ and $
	\tilde{\theta}_t = \hat{\theta}_t$, the center of the confidence ellipsoid, for all $t\in[T]$.
\end{example}

Before describing how TS can be derived as an instance of \POFUL{}, we introduce a definition.
%
\begin{definition}
Let $\delta \in (0,1)$. We define $\cD^{SA}(\delta)$ as a distribution satisfying 
	$	 
	\PP_{\eta \sim \cD^{SA}(\delta)}\left[\|\eta\| \leq 1\right] \geq 1 - \delta
	$.
\end{definition}


\begin{example}[TS]\label{ex:ts}	Linear Thompson Sampling (LinTS) algorithm is a generic randomized algorithm that samples from a distribution constructed from the RLS estimate at each step. At time $t$, LinTS samples as 
%
$
\tilde{\theta}_{t}  =\hat{\theta}_t +\iota^{TS}_t \beta_{t,\delta^\prime, \lambdareg}^{RLS}  V_t^{-1/2} \eta_t
$,
%
where $\delta^\prime = \delta/2T$,  $\iota^{TS}_t$ is inflation parameter controlling the scale of the sampling range, and $\eta_t$ is a random sample from a normalized sampling distribution $\cD^{SA}(\delta^\prime)$ that concentrates with high probability. LinTS is a special case of \POFUL{} where  $\iota_t = \iota^{TS}_t$, $\tau_t = 0$ and 
$
\tilde{\theta_t} = \hat{\theta}_t +\iota^{TS}_t \beta_{t,\delta^\prime, \lambdareg}^{RLS}  V_t^{-1/2} \eta_t
$.
%
Setting  $\iota_t = \tilde{\cO}(1)$ corresponds to the original LinTS algorithm, while  setting $\iota_t= \tilde{\cO}(\sqrt{d})$  corresponds to the frequentist variant of LinTS studied in \cite{agrawal2013thompson,abeille2017linear}, namely TS-Freq. This means TS-Freq inflates the posterior by a factor of order $\sqrt{d}$.
\end{example}

\begin{example}[Greedy]\label{ex:greedy}
Greedy is a special case of \POFUL{} with $\iota_t = \tau_t = 0$, $\tilde{\theta}_t = \hat{\theta}_t$,  $\forall t$.
\end{example}


























% todo add a remark that this is equivalent to adding a random vector


